% Part: incompleteness
% Chapter: arithmetization-syntax
% Section: introduction

\documentclass[../../../include/open-logic-section]{subfiles}

\begin{document}

% SEGMENT OLP-0281-S01

\olfileid{inc}{art}{int}
\olsection{அறிமுகம்}

கணிப்புத்தன்மையையும் தருக்கத்தையும் இணைக்க, தருக்கத்தின் பொருள்கள் (குறிகள்,
சொற்கள், !!{formula}s, !!{derivation}s), அவற்றின் மீதான செயல்கள்,
அவற்றின் பண்புகள், தொடர்புகள் ஆகியவற்றைப் பற்றிக் கணக்கீட்டு முறையில்
கையாளத்தக்கவாறு பேசும் வழி நமக்குத் தேவை. குறிகள், குறிகளின் தொடர்கள்,
அவற்றிலிருந்து கட்டப்பட்ட பிற பொருள்கள் ஆகியவற்றின் மீதான கணிக்கத்தக்க
சார்புகளையும் தொடர்புகளையும் கருதுவதன் மூலம் இதை நேரடியாகச் செய்யலாம்.
தருக்கத் தொடரமைப்பின் பொருள்கள் அனைத்தும் முடிவுறுபவையாகவும், குறிகளின்
!!a{enumerable} கணங்களிலிருந்து கட்டப்படுபவையாகவும் இருப்பதால், சில
கணக்கீட்டு மாதிரிகளில் இது சாத்தியம். ஆனால் மீள்வரையறைச் சார்புகள் போன்ற
பிற கணக்கீட்டு மாதிரிகள் எண்களுக்கும் அவற்றின் தொடர்புகளுக்கும் சார்புகளுக்கும்
வரம்பிடப்பட்டுள்ளன. மேலும், இறுதியில் சில கோட்பாடுகளுக்குள்ளும், குறிப்பாக
எண்கணித மொழியில் உருவாக்கப்பட்ட கோட்பாடுகளுக்குள்ளும் தொடரமைப்பைக் கையாள
விரும்புகிறோம். இந்நேர்வுகளில் தொடரமைப்பை \emph{எண்கணிதமாக்குவது} தேவை;
அதாவது தொடரமைப்புப் பொருள்கள், அவற்றின் மீதான செயல்கள், அவற்றின் தொடர்புகள்
ஆகியவற்றை முறையே எண்கள், எண்கணிதச் சார்புகள், எண்கணிதத் தொடர்புகள் எனப்
பிரதிநிதித்துவப்படுத்த வேண்டும். லைப்னிட்ஸிடம் தொடங்கும் இவ்வெண்ணம்,
தொடரமைப்புப் பொருள்களுக்கு எண்களை ஒதுக்குவதாகும்.

குறிகளுக்கு அவற்றின் ``குறியீடுகளாக'' எண்களை ஒதுக்குவது ஒப்பீட்டளவில்
நேரடியானது. சில குறிகள் சிறிது சவாலாக அமைகின்றன; ஏனெனில், எடுத்துக்காட்டாக,
முடிவிலியாகப் பல !!{variable}s உள்ளன; ஒவ்வோர் உறுப்பெண்~$n$ க்கும்
முடிவிலியாகப் பல !!{function}s கூட உள்ளன. ஆனால் $\Obj v_2$ க்கும்
$\Obj v_3$ க்கும் வெவ்வேறு குறியீடுகள் ஒதுக்கப்படுமாறு, குறிகளுக்கு எண்களை
முறைப்படி ஒதுக்குவது இயல்பாகவே சாத்தியம். குறிகளின் தொடர்கள் (சொற்களும்
!!{formula}s உம் போன்றவை) இன்னும் பெரிய சவால். ஆனால் எண் தொடர்களை
முற்றிலும் எண்கணிதமுறையில் கையாள முடிந்தால் (எடுத்துக்காட்டாக, தொடர்களின்
பகா எண் அடுக்குக் குறியீடாக்கம் மூலம்), தனிக் குறிகளின் குறியீடாக்கத்தை
குறிகளின் தொடர்களுக்கான குறியீடாக்கமாகவும், அதன்பின் !!{formula}s இன்
தொடர்கள் அல்லது !!{derivation}s போன்ற பிற அமைப்புகளுக்கான குறியீடாக்கமாகவும்
விரிவாக்கலாம். இவ்விரிவாக்கப்பட்ட குறியீடாக்கம் ``கோடல் எண்ணிடல்''
எனப்படுகிறது. ஒவ்வொரு சொல்லுக்கும், !!{formula} க்கும்,
!!{derivation} க்கும் ஒரு கோடல் எண் ஒதுக்கப்படுகிறது.

குறிகளின் தொடர்களை அவற்றின் குறியீடுகளின் தொடர்களாகக் குறியீடாக்கி,
கணிக்கத்தக்க சார்புகளால் கையாளக்கூடிய ஒரு தொடர்-குறியீடாக்க முறைமையைத்
தேர்ந்தெடுத்தால், கணிக்கத்தக்க சார்புகளைக் கொண்டு கோடல் எண்களையும் கையாளலாம்.
நடைமுறையில், தொடர்புடைய எல்லாச் சார்புகளும் முதன்மை மீள்வரையறைச் சார்புகளாக
இருக்கும். எடுத்துக்காட்டாக, ஒரு தொடரின் நீளத்தைக் கணிப்பதும், தொடரின்
குறியீட்டிலிருந்து அதன் $i$-ஆம் உறுப்பைக் கணிப்பதும் முதன்மை மீள்வரையறைச்
சார்புகள். தொடரைக் குறியீடாக்கும் எண், எடுத்துக்காட்டாக, !!a{formula}~$!A$
இன் கோடல் எண் எனில், ஒரு !!a{formula} வின் நீளத்தையும், அந்த
!!a{formula} விலுள்ள $i$-ஆம்
குறியின் (குறியீட்டையும்) $!A$ இன் கோடல் எண்ணிலிருந்து கணிக்கலாம் என்பதை
உடனே காண்கிறோம். ஒரு எண் சரியாக உருவாக்கப்பட்ட சொல்லின் கோடல் எண்ணாகவோ,
சரியான !!{derivation} இன் கோடல் எண்ணாகவோ இருக்கும் பண்பு முதன்மை
மீள்வரையறையானது என்று நிறுவுவது சற்றுக் கடினம். இருப்பினும் இது சாத்தியம்;
ஏனெனில் நமக்கு வேண்டிய தொடர்கள் (சொற்கள், !!{formula}s,
!!{derivation}s) தொகுத்தறிதல்வழி வரையறுக்கப்பட்டுள்ளன.

ஓர் எடுத்துக்காட்டாக இடைநிறுத்தல் செயலைக் கருதுக. $!A$ ஒரு வாய்பாடாகவும்,
$x$ ஒரு மாறியாகவும், $t$ ஒரு சொல்லாகவும் இருந்தால், $!A$ இல்~$x$ கட்டற்று
வரும் ஒவ்வொரு இடத்திலும்~$t$ ஐ இடுவதால் கிடைப்பதே $\Subst{!A}{t}{x}$.
இப்போது $!A$, $x$, $t$ ஆகியவற்றுக்கு முறையே $k$, $l$, $m$ என்ற கோடல்
எண்களை ஒதுக்கியுள்ளோம் என்க. அதே திட்டம் $\Subst{!A}{t}{x}$ க்கு,
எடுத்துக்காட்டாக~$n$ என்ற கோடல் எண்ணை ஒதுக்குகிறது. $k$, $l$, $m$
ஆகியவற்றை~$n$ குக் கொண்டுசெல்லும் இந்தப் பொருத்தம், இடைநிறுத்தல் செயலின்
எண்கணித ஒப்புமையாகும். இடைநிறுத்தல் செயல் $!A$, $x$, $t$ ஆகியவற்றை
$\Subst{!A}{t}{x}$ குக் கொண்டுசெல்லும்போது, எண்கணிதமாக்கப்பட்ட இடைநிறுத்தல்
சார்பு, $k$, $l$, $m$ என்ற கோடல் எண்களை~$n$ என்ற கோடல் எண்ணுக்குக்
கொண்டுசெல்கிறது. இச்சார்பு முதன்மை மீள்வரையறையானது என்று காண்போம்.

தொடரமைப்பின் எண்கணிதமாக்கம் அருவமான ஆர்வத்துக்கு மட்டும் உரியதல்ல. மொழிகளைப்
``பற்றிப் பேசுவதற்கான'' பொறிமுறைகள் இல்லாத எண்கணித மொழி போன்ற மொழிகளால்கூட
கூற்றுகளின் சிக்கலான பண்புகளை முறைப்படுத்த முடியும் என்பது தொடக்கத்தில்
எளிதல்லாத நுண்ணறிவாக இருந்தது. பின்னர், பியானோ எண்கணிதம் போன்ற இம்மொழியிலுள்ள
ஒரு கோட்பாடு அதன் சொந்த மொழியைப் பற்றி (எடுத்துக்காட்டாக, !!{sentence}s
வருவிக்கத்தக்கவையா அல்லது மெய்யா என்பது உட்பட) எதை \emph{நிறுவ} முடியும்
என்று கேட்பது ஒரு சிறு படியே. இது கோடலின் (வருவிக்கவியலாமை பற்றிய) புகழ்பெற்ற
வரம்புத் தேற்றங்களுக்கும், டார்ஸ்கியின் (மெய்மையை வரையறுக்கவியலாமை பற்றிய)
புகழ்பெற்ற வரம்புத் தேற்றத்துக்கும் நம்மைக் கொண்டுசெல்கிறது. ஆனால் தொடரமைப்பை
எண்கணிதமாக்கும் உத்தி, கோட்பாடுகளின் கணக்கீட்டுத் திறன் அல்லது வெவ்வேறு
கணிப்புத்தன்மை மாதிரிகளுக்கிடையிலான உறவு போன்ற கணிப்புத்தன்மைக் கோட்பாட்டின்
முக்கிய முடிவுகளை நிறுவுவதற்கும் தேவை. பிற பொருள்களையும் பண்புகளையும்
எண்கணிதமாக்குவதற்கான மாதிரியாகத் தொடரமைப்பின் எண்கணிதமாக்கம் அமைகிறது.
எடுத்துக்காட்டாக, அமைவுநிலைகளையும் கணக்கீடுகளையும் (டியூரிங் பொறிகளின்
அமைவுநிலைகளையும் கணக்கீடுகளையும் போன்றவை) அதேபோல் எண்கணிதமாக்க முடியும்.
இது ஒரு மாதிரியிலுள்ள (எடுத்துக்காட்டாக டியூரிங் பொறிகள்) கணக்கீடுகளை மற்றொரு
மாதிரியில் (எடுத்துக்காட்டாக மீள்வரையறைச் சார்புகள்) உருவகப்படுத்த உதவுகிறது.

\end{document}
